\begin{statement}{A.5.11}{Proved}{The genus of a horizontal conductor image from a local torus group}
Let $f:Y\to C_0$ be a proper holomorphic map from a smooth
compact complex threefold to a smooth curve, with connected
fibres and complex two-torus general fibre. Let $S$ be an
irreducible reduced fibre support, $\eta:T\to S$ its
normalization, and $\pi:Z\to T$ a resolution which is
irrational ruled, with ruling $p:Z\to B$.
Suppose every exceptional curve of $\pi$ is vertical, and
the strict conductor has just one horizontal irreducible
component $C$, which is a section of $p$.
Let $D=\eta\pi(C)$ and let $\Sigma=D^\nu$ be its normalization.
Then $g(\Sigma)\le1$.
In particular, when $C\to D$ has degree one, $B$ is elliptic.
This applies to the horizontal-cusp case of $D_Z=2C+V$;
it allows arbitrary multiplicity of the full fibre and
arbitrary singularities or nonreduced conductor image on
the finitely many vertical curves.
\end{statement}
\begin{proof}
There is nothing to prove when $\Sigma$ is rational, so
assume $g(\Sigma)\ge1$.
Statement A.5.8 gives a holomorphic ruling
$\widetilde p:T\to B$ with $p=\widetilde p\pi$.
Identifying $C$ with $B$ by its section map, normalization
of its nonconstant map to $D$ gives a finite surjective
holomorphic map $\phi:B\to\Sigma$.
Put $\psi=\phi\widetilde p:T\to\Sigma$.

We describe precisely the finite identifications needed to
make this map descend to $S$. Every vertical conductor
component of $T$ lies over one of finitely many points of
$B$, so $\psi$ is constant on it. Its image under $\eta$
has rational normalization: a vertical component is the
image of a rational curve on $Z$, and a nonconstant map
from that curve to its image normalization is finite.
Such an image cannot equal $D$, because $g(\Sigma)\ge1$.
Consequently these vertical images meet $D$ at finitely
many points. Away from those points and from the finitely
many singular points of $D$, every normalization preimage
of a point of $D$ lies on $C$, and its $\psi$ value is the
unique corresponding point of $\Sigma$.
At a point outside $D$ with more than one normalization
preimage, all those preimages lie on vertical conductor
components. There are no other horizontal conductor
components by hypothesis. Thus there is a finite set
$A\subset\Sigma$ such that, for every $s\in S$, either
$\psi(\eta^{-1}(s))$ is a singleton or it is contained
in $A$. One can take $A$ to contain the constant values
on all vertical conductor components and all points of
$\Sigma$ above the singularities of $D$ and its
intersections with the vertical image curves.

Choose an embedded finite tree $J$ in the underlying
oriented surface $\Sigma$ containing $A$, and let
$q:\Sigma\to\Sigma/J$ collapse that tree to one point.
Such a tree is obtained by joining the finite set by
arcs and retaining a spanning tree after subdividing
their finitely many crossings. Triangulate the surface
with $J$ a subcomplex. A contractible subcomplex is a
cofibration; contracting it gives a homotopy equivalence
of this CW pair. Concretely the homotopy extension
property extends a contraction of $J$ and supplies a
homotopy inverse to the quotient map. In particular
\[
              q_*:\pi_1(\Sigma)\xrightarrow{\sim}
                         \pi_1(\Sigma/J).          \tag{A.Pi.1}
\]
The map $q\psi$ is constant on every fibre of the finite
proper surjection $\eta$. A proper continuous surjection
between these compact Hausdorff spaces is a quotient
map, so there is a unique continuous map
\[
 a:S\longrightarrow\Sigma/J,
 \qquad a\eta=q\psi.                              \tag{A.Pi.2}
\]
Let $\delta:\Sigma\to D\subset S$ be the curve normalization
map. On its dense smooth locus, lifting through $C$ gives
$a\delta=q$. Continuity gives this equality everywhere.
Equation (A.Pi.1) therefore makes $a_*$ surjective on
fundamental groups, without any assertion that a finite
point-identification quotient has the same fundamental
group as $\Sigma$.

We next extend the actual map (A.Pi.2) to a neighbourhood
of $S$ in $Y$. The simultaneous analytic triangulation
theorem of \cite[\S3, Theorem~2, p.~464]{Lojasiewicz1964}
triangulates the pair $(Y,S)$, since a complex analytic
subset is real semianalytic. Compactness makes the
subcomplex representing $S$ finite. After barycentric
subdivision it is full: vertices lying in it span a
face in that same subcomplex. Write $\lambda_v$ for
barycentric coordinates and let $L$ be its vertex set.
On the open neighbourhood where $\sum_{v\in L}\lambda_v>0$,
the formula
\[
 r(x)=\frac{\sum_{v\in L}\lambda_v(x)v}
                  {\sum_{v\in L}\lambda_v(x)}       \tag{A.Pi.3}
\]
defines a continuous retraction to $S$ in triangulation
coordinates. Its value lies in the face spanned by the
$L$ vertices of the same simplex, so these formulas
agree across faces. This supplies an actual open
neighbourhood $N$ and a map $r:N\to S$ fixed on $S$.

Let $c=f(S)$. Properness allows a small disc $\Delta$
about $c$ with $U=f^{-1}(\Delta)\subset N$ and no other
critical value. The map $(ar)|_U:U\to\Sigma/J$ is
surjective on fundamental groups: its restriction to
$\delta(\Sigma)\subset S\subset U$ is $q$, whose
fundamental-group map is an isomorphism.
Set $U^*=U\setminus S=f^{-1}(\Delta^*)$.
The inclusion induces a surjection
$\pi_1(U^*)\twoheadrightarrow\pi_1(U)$.
Indeed a loop in the smooth six-manifold $U$ can be
perturbed relative to its base point to be transverse
to an analytic stratification of $S$; every stratum
has real codimension at least two, so the resulting
one-dimensional loop misses every stratum.

Over $\Delta^*$, the proper submersion is a smooth
torus bundle. The fundamental-group exact sequence
and $\pi_2(\Delta^*)=0$ give
\[
 1\longrightarrow\Z^4\longrightarrow\pi_1(U^*)
       \longrightarrow\Z\longrightarrow1.
\]
A lift of the generator of $\Z$ splits this group
extension, hence
\[
 \pi_1(U^*)\simeq\Z^4\rtimes_T\Z
 \twoheadrightarrow\pi_1(U)
 \twoheadrightarrow\pi_1(\Sigma/J)
 \simeq\pi_1(\Sigma).                             \tag{A.Pi.4}
\]
The first group is metabelian: its commutator subgroup
is contained in the abelian normal subgroup $\Z^4$.
Every quotient in (A.Pi.4) is therefore metabelian.
A closed orientable surface of genus at least two
has a fundamental group surjecting onto the free
group on two generators: in its usual surface
presentation, retain two $a$ generators and send all
$b$ generators and the other $a$ generators to the
identity. That free group surjects onto $S_4$ by
sending its generators to $(1234)$ and $(12)$;
their conjugates supply the adjacent transpositions.
The commutator subgroup of $S_4$ is the nonabelian
group $A_4$, so this quotient is not metabelian.
This contradicts (A.Pi.4). Thus $g(\Sigma)=1$ in the
positive-genus case. When $C\to D$ has degree one,
$\Sigma\simeq C\simeq B$, which proves the final
assertion.
\end{proof}


% A.5.12 from part_A.tex

